\chapter{Introduction}
\section{Context}
\subsection{Mixed-Integer Linear Programming}
Linear optimization \cite{wiki:Linear_programming} is a method in which a mathematical model formulating a problem is optimized given requirements, constraints, and an objective function in linear form. It is heavily used in mathematics and computer science, but its applications can also be actively seen in real life \cite{kunwar2022introduction}, such as in agriculture for production planning, in industries to maximize efficiency, and in business economics to maximize profit. These are some examples where LP (Linear Programming) is used. 

When we talk about linear optimization, it is usually continuous, but not all decisions are continuous in nature. Integer Linear Programming \cite{wiki:Integer_programming} (ILP) includes integer variables which allows for concrete solutions. A mix of both is called Mixed-Integer Linear Programming (MILP), with some discrete and some continuous variables.

\subsection{Optimization Platform}
Solving such problems mathematically requires expertise, and as we have seen, MILP is widely used in real life; therefore, a more accessible approach to solving such models is warranted. This is where the Optimization Platform \cite{MarkoThesisMsc, DeniThesisMsc, MarkoThesisBsc, MaksimThesisMsc} comes in. It is a web-based platform used to model and solve MILPs. It provides many features that make the modeling process easier and more accessible, such as a drag-and-drop modeler, the ability to create multiple workspaces to explore different models simultaneously, and the possibility of multi-user collaboration. All of these features share the same goal: making the user’s work simpler and more efficient.

\subsection{Human-AI Collaboration}
However, recent advancements in artificial intelligence, machine learning and the development of large language models \cite{vaswani2023attentionneed} have opened the door to promising, previously unseen possibilities for human–machine collaboration. Following our goal of making this platform more accessible and productive for users, we explore the possibilities of closing the gap between artificial intelligence and the optimization modeler.

\section{Challenges}

The current feature set of the optimization platform has already made it very accessible and easy to use. However, there are still several aspects that can be improved to take the platform to the next level. Below are some of the existing challenges that this work aims to address:

\begin{enumerate}
    \item \textbf{Inefficiency:} As the primary method of modeling is via the drag-and-drop modeler interface, relying solely on it for large and complex models can be very time-consuming and redundant. Although this approach makes modeling more intuitive, it significantly decreases efficiency. Increased model complexity also raises the likelihood of errors due to limited supervision and cross-verification.
    
    \item \textbf{Inaccessibility:} With an increasing number of features on the platform, it becomes difficult for users to familiarize themselves with all available functionalities before they can start building a model. This results in a steep learning curve and a high entry barrier, thereby contributing to the inaccessibility of the platform.
    
    \item \textbf{Lack of Support:} While human collaborators can join a workspace and work together in teams to support each other, users working independently have no immediate assistance to rely on when they encounter difficulties. Providing live support for such situations would have a positive impact on multiple aspects of the user experience and significantly improve overall usability.
\end{enumerate}


\section{Goal}

The goal of this work is to integrate an AI agent into the optimization platform so that it can act as a first-class collaborator alongside other users in a shared workspace. The AI agent can be asked, in natural language, to make changes that users want to apply to the current state of the optimization model. The agent then processes the user’s request, generates a solution, and applies it to the shared workspace.

After the AI agent applies the changes, they become visible to the users but remain tentative, pending user feedback or approval. Users can review the proposed solution generated by the AI agent and either accept it, reject it, or provide feedback on the generation. This architecture preserves collaboration, human oversight, supervision, and validation.

\section{Solution Outline}

This section provides a high-level overview of the proposed solution and explains how its individual components work together. We first describe the system flow at an abstract level and then examine each component in greater detail in the subsequent sections.

\subsection{AI-Agent Collaborator}
A user can begin using the Optimization Modeler by logging in as a master user and creating an optimization model. At this stage, the user can invite an AI agent to the workspace through the same collaboration layer that is already available for human collaborators. The AI agent joins the shared modeling workspace in the same manner as a human user; however, unlike a human collaborator, the AI agent can independently process requests and make modifications to the shared workspace.

\subsection{Iterative Feedback}
The AI agent is powered by a large language model that interprets user input expressed in natural language and translates it into corresponding changes in the workspace. The user retains full control over this process: when the user issues a request, the AI agent proposes a solution that is displayed in the workspace, after which the user can decide whether to accept or reject the generated changes. The user may also provide additional feedback, which the AI agent uses to iteratively refine and improve its output. This architecture ensures user oversight and control while maintaining a state-of-the-art, stateful, agentic system design.

\subsection{Safety and Validation}
In addition, a validation layer is applied after the AI model generates its output and before the changes are committed to the workspace. This layer validates the generated model for syntactic and structural correctness, allowing invalid or hallucinated outputs to be discarded. As a result, the validation layer adds an additional safeguard to the overall architecture.



